\section{Preliminaries}\label{sec:prelim}
This section summarises the concepts and methods on which PHQD is built, so that the later sections can be read without prior knowledge of quantum computing.

A convolutional autoencoder consists of an encoder $E$ and a decoder $D$ trained to reconstruct images through a compact bottleneck~\cite{rombach}. In our model $E$ maps a $1\times28\times28$ image $x$ to a latent map $z=E(x)\in\mathbb{R}^{64\times7\times7}$ using two strided convolutions with batch normalisation and LeakyReLU activations, and $D$ maps $z$ back to image space. The latent map is corrupted with additive Gaussian noise, $\tilde z=z+\sigma\epsilon$ with $\epsilon\sim\mathcal{N}(0,I)$ and $\sigma=0.3$, which is the same kind of corruption that diffusion models learn to remove~\cite{ho}. A denoiser $f_\theta$ is trained to predict $z$ from $\tilde z$ by minimising the mean squared error (MSE) $\lVert f_\theta(\tilde z)-z\rVert_2^2$, and the denoised latent map is passed through $D$ to obtain the final image.

A latent map with 3,136 values is too large to encode into a small circuit at once. We therefore slide a $2\times2$ window with stride one over the $7\times7$ spatial grid, which gives $6\times6=36$ overlapping patches, each holding $64\times2\times2=256$ values. Every patch is denoised separately, and values that fall in several patches are averaged when the latent map is reassembled. The overlap lets neighbouring patches share information and reduces visible seams.

A qubit has a state $\alpha\lvert0\rangle+\beta\lvert1\rangle$ with $\lvert\alpha\rvert^2+\lvert\beta\rvert^2=1$, and a register of $n$ qubits is changed by unitary gates. The single-qubit rotation $R_Y(\theta)=e^{-i\theta Y/2}$ turns a real number $\theta$ into a state change, and the general rotation $\mathrm{Rot}(\phi,\theta,\omega)=R_Z(\omega)R_Y(\theta)R_Z(\phi)$ has three trainable angles. The controlled-NOT (CNOT) gate flips a target qubit when a control qubit is $\lvert1\rangle$ and creates entanglement; we apply it in a ring over the qubits. Measuring a Pauli operator $P\in\{X,Y,Z\}$ on qubit $i$ yields the expectation value $\langle P_i\rangle\in[-1,1]$, and these bounded real numbers are the outputs passed back to the classical network.

A PQC applies a data-dependent encoding followed by trainable layers, $U(x;w)$, and returns $f_i(x;w)=\langle0\rvert U^\dagger(x;w)\,P_i\,U(x;w)\lvert0\rangle$~\cite{cerezo,benedetti}. Its trainable weights are updated by gradient descent in the same way as network weights, using either the parameter-shift rule or back-propagation through a classical state-vector simulator. All circuits in this work are simulated classically, so the effects of hardware noise on NISQ devices~\cite{preskill} are not included. Random deep circuits may have vanishing gradients~\cite{mcclean}, which is one reason we keep the circuits shallow.

In angle encoding, each classical feature sets the rotation angle of one qubit, so $n$ features need $n$ qubits. The IQP-style encoding of Havl\'{\i}\v{c}ek et al.~\cite{havlicek} applies Hadamard gates, feature-dependent $R_Z$ rotations and two-qubit $ZZ$ interactions that depend on products of features, which adds nonlinear feature interactions. In data re-uploading~\cite{salinas}, the same features are encoded again in every layer between trainable gates. This enlarges the set of frequencies that the model can represent~\cite{schuld}. Reading out $X$, $Y$ and $Z$ on every qubit instead of $Z$ only increases the number of output features from $n$ to $3n$ for the same circuit.

In our hybrid denoiser a classical network first projects each 256-value patch to $n$ rotation angles. The PQC then maps these angles to expectation values, and a second classical network, together with a linear residual path, maps the circuit output back to a 256-value patch. A stitching network, a small convolutional network, then refines the reassembled latent map, and in the 6-qubit model it is trained jointly with the denoiser.

The multilayer perceptron (MLP) denoises patches with fully connected layers. The U-Net~\cite{unet} is an encoder--decoder convolutional network with skip connections that denoises the whole latent map, and linear regression gives the simplest linear mapping from a noisy patch to a clean one. Because the classical layers of the hybrid model have far more parameters than the circuit, these baselines show how much of its accuracy a purely classical model can reach.

Reconstruction quality is measured by MSE, by the peak signal-to-noise ratio $\mathrm{PSNR}=10\log_{10}(\mathrm{MAX}^2/\mathrm{MSE})$ and by the structural similarity index (SSIM), which compares local luminance, contrast and structure. The Fr\'echet Inception Distance (FID) compares the mean and covariance of Inception-V3 features of generated and real images, where lower values mean that the two sets are more alike. Higher PSNR and SSIM, and lower MSE and FID, indicate better results.
